\section{Tree Prikry forcing in the native order convention}
\label{sec:p06-tree-prikry-forcing}

The tree presentation allows the measure-one set of possible next points
to depend on the finite stem. It therefore supplies direct decision without
assuming a normal ultrafilter. This is Gitik's native tree-Prikry route,
\cite[\S1.2, Definitions~1.14--1.17, Lemmas~1.18--1.24 and
Theorem~1.25]{GitikPrikryAuthorDraft}. The source's \(\geq,\geq^*\)
orders are retained here. A separately named host order transports the
existing recursive forcing machinery; it does not replace those glyphs.

\begin{hbconvention}{def:p06-tree-prikry-ground}{Nonprincipal completeness, not normality}
Let \(M\) be a transitive set model of \(\mathsf{ZFC}\) as in
\cref{def:p06-ftr-ground}. In \(M\), let \(\kappa\) be an uncountable
cardinal and \(U\) a nonprincipal \(\kappa\)-complete ultrafilter on
\(\kappa\). Normality is not a standing hypothesis. Section~1.2 drops
normality, not the nonprincipality of the preceding measure setup;
nonprincipality is needed for unbounded successors. All tree families
and completeness assertions below are ground computations. When an
extension is used, supply an ambient \(M\)-generic filter. External
countability can produce one but is not an internal proof hypothesis.
No such ground or measure is inferred from \(\operatorname{Con}(\mathsf{ZFC})\).
\end{hbconvention}

\subsection{Native trees, cones, and orders}

\begin{hbdefinition}{def:p06-tree-prikry-trees}{The source U-tree and its trunk}
A \(U\)-tree \(T\) with trunk \(t\) consists of finite strictly
increasing sequences below \(\kappa\), ordered by end extension.
The trunk belongs to \(T\) and is comparable by end extension with
every node. Above the trunk, a tree retains the intermediate initial
segments of every node. For every node \(\eta\) end-extending \(t\),
\[
 \operatorname{Suc}_T(\eta)=
   \{\alpha<\kappa:\eta\mathbin{\frown}\langle\alpha\rangle\in T\}\in U.
\]
Write \(\operatorname{Lev}_n(T)=\{\eta\in T:
\operatorname{length}(\eta)=n\}\), retaining the source's absolute
level indexing. Its finite-stem notation \(|t|\) denotes the length,
which also equals the cardinality of the finite sequence as a function.
For \(r\in T\) end-extending the trunk, the source cone is
\[
 (T)_r=\{r'\in T:r'\text{ end-extends }r\}.
\]
It is a rooted tree with trunk \(r\); ancestors below \(r\) need not
be retained in a cone. Thus rooted prefix closure above a trunk is not
confused with compulsory global closure down to the empty sequence.
\end{hbdefinition}

\begin{hbdefinition}{def:p06-tree-prikry-native-orders}{Greater is stronger in both native orders}
The source condition set \(P\) consists of pairs \(\langle t,T\rangle\)
where \(T\) is a \(U\)-tree with trunk \(t\). For
\(p=\langle t,T\rangle\), \(q=\langle s,S\rangle\), put
\[
 p\geq q\quad\Longleftrightarrow\quad T\subseteq S,
 \qquad
 p\geq^*q\quad\Longleftrightarrow\quad T\subseteq S\text{ and }t=s.
\]
The inverse relations are \(\leq,\leq^*\), as in the source.
The full order allows a longer trunk; the direct order fixes it.
The weakest condition is the empty trunk with the tree of all finite
increasing sequences in \(\kappa\). The condition collection is a
ground set, by Separation over the finite-sequence set times its power
set. A native dense set has a member
\(q\geq p\) above every \(p\); a native open set stays closed under
stronger \(q\geq p\). A native generic filter is directed toward
common stronger conditions and closed toward weaker conditions.
\end{hbdefinition}

\begin{hblemma}{thm:p06-tree-prikry-structure}{proved=annex:p06.tree-prikry.structure}{Trunk extension, intersection, grafting and the chain condition}
Each rooted cone is again a \(U\)-tree. The two native relations are partial orders, and
\(\geq^*\subseteq\geq\). If \(\langle t,T\rangle\geq\langle s,S\rangle\),
then \(t\) end-extends \(s\) and belongs to \(S\).
An intersection of fewer than \(\kappa\) same-trunk \(U\)-trees is
a same-trunk \(U\)-tree. Replacing independently each cone at one
fixed finite level by a direct \(U\)-subtree gives a direct tree extension.
The direct order has bounds for strengthening chains of length less than
\(\kappa\), and the full forcing has \(\kappa^+\)-cc.
The measure is uniform, and \(\kappa\) is regular and a strong limit,
without using normality.
\end{hblemma}

\begin{hblemma}{thm:p06-tree-prikry-host}{proved=annex:p06.tree-prikry.host}{An explicit source-to-host order-dual bridge}
On the same condition set define
\(q\leq_h p\) exactly when \(q\geq p\), and
\(q\leq_h^*p\) exactly when \(q\geq^*p\).
Compatibility, dense sets and generic filters are identical under these
two stated orientations. Ground names and their evaluation are unchanged.
The native recursive forcing relation, with strengthening quantifiers
\(q\geq p\), agrees with the host relation with quantifiers
\(q\leq_h p\). Hence the host density, truth and ZFC-transfer results
apply to this native presentation. There is no implicit forcing with the
auxiliary direct order.
\end{hblemma}

\subsection{Level-by-level thinning without diagonal normality}

\begin{hblemma}{thm:p06-tree-prikry-level-colors}{proved=annex:p06.tree-prikry.level-colors}{A homogeneous U-subtree at every finite level}
If \(f:T\to\lambda\), \(\lambda<\kappa\), there is a direct
subtree \(S\subseteq T\) on which \(f\) is constant on each
\(\operatorname{Lev}_{|t|+n}(S)\), \(n<\omega\).
Different levels may have different colors. The construction propagates
each finite target level's colors backward through ultrafilter-large
successor fibers, then intersects countably many same-trunk trees.
It does not assert the classical global-reservoir homogeneity lemma for
a non-normal measure.
\end{hblemma}

\begin{hbtheorem}{thm:p06-tree-prikry-decision}{proved=annex:p06.tree-prikry.decision}{The native tree Prikry condition}
For a fixed forcing sentence with ground names and any
\(\langle t,T\rangle\), there is
\(\langle t,S\rangle\geq^*\langle t,T\rangle\) deciding it in the
full forcing. The proof constructs the source's level-by-level sequence
of trees, freezes previously reached levels, synchronizes decisions at
siblings, and eliminates a minimal undecided trunk length. Every
countable intersection is taken pointwise at each successor set.
\end{hbtheorem}

\begin{hblemma}{thm:p06-tree-prikry-strong-dense}{proved=annex:p06.tree-prikry.strong-dense}{A uniform cone level in a native dense open set}
For a native dense open \(D\) and \(\langle t,T\rangle\), there
are a direct subtree \(S\) and \(n<\omega\) such that for every
\(m\geq n\) and every
\(\eta\in\operatorname{Lev}_{|t|+m}(S)\),
\[
 \langle\eta,(S)_\eta\rangle\in D.
\]
The proof gives the cone witnesses and level grafts explicitly. It does
not infer this strong statement by citing the preceding decision theorem.
\end{hblemma}

\subsection{Branches and preservation}

\begin{hbtheorem}{thm:p06-tree-prikry-branch}{proved=annex:p06.tree-prikry.branch}{The generic branch and its exact filter}
The union \(C\) of generic trunks is an increasing \(\omega\)-sequence
cofinal in \(\kappa\), new over \(M\), and eventually in every
ground \(A\in U\). The generic filter is exactly
\[
 G(C)=\{\langle t,T\rangle:t\text{ is initial in }C,
       \ (\forall n\geq|t|)\ C\mathbin{\upharpoonright}n\in T\}.
\]
Thus \(M[G]=M[C]\). Here the branch condition is membership in a
whole succession tree, not containment in one fixed measure-one tail.
\end{hbtheorem}

\begin{hbtheorem}{thm:p06-tree-prikry-preservation}{proved=annex:p06.tree-prikry.preservation}{No bounded subsets and no cardinal collapse}
For every ground \(\gamma<\kappa\),
\(\mathcal P(\gamma)^{M[G]}=\mathcal P(\gamma)^M\).
All ground cardinals are preserved, and for each ground limit ordinal
\(\delta\) its cofinality changes to \(\omega\) exactly when its
ground cofinality was \(\kappa\); otherwise its cofinality is unchanged.
This is proved for the stated nonprincipal \(\kappa\)-complete measure,
not imported from the normal-measure classical theorem.
\end{hbtheorem}

\begin{hbtheorem}{thm:p06-tree-prikry-family-criterion}{proved=annex:p06.tree-prikry.family-criterion}{A stem-indexed tree-family genericity criterion}
For an ambient increasing \(\omega\)-sequence \(C\subseteq\kappa\),
the displayed \(G(C)\) is generic exactly when the following holds:
for every ground assignment \(s\mapsto T_s\) of a \(U\)-tree with
trunk \(s\) to every finite increasing \(s\), some initial stem
\(s=C\mathbin{\upharpoonright}N\) has
\(C\mathbin{\upharpoonright}n\in T_s\) for every \(n\geq N\).
This criterion is a local consequence of the tree construction. The
necessary eventual-containment condition for individual \(A\in U\)
is not silently substituted for it when normality is absent.
\end{hbtheorem}

\begin{hbtheorem}{thm:p06-tree-prikry-normal-comparison}{proved=annex:p06.tree-prikry.normal-comparison}{The additional normal-measure comparison}
Assume now, additionally, that \(U\) is normal. The map taking a
classical host condition \((s,A)\) of \cref{def:p06-prikry-poset}
to the tree with trunk \(s\) and all
later points in \(A\) is a dense embedding: a classical host extension
\(q\leq p\) corresponds to the native tree extension
\(\iota(q)\geq\iota(p)\). It preserves and reflects incompatibility,
and its generic extensions are the same, with the same cofinal sequence.
Via \cref{thm:p06-prikry-order-bridge}, the source finite-set conditions
of \cref{def:p06-prikry-gitik-orders} correspond as well, preserving
their native \(\geq,\geq^*\) orientations rather than the host glyphs.
Under this added normality hypothesis, eventual containment in every
\(A\in U\) implies the tree-family criterion. This recovers the
Mathias comparison at its actual normal-measure premise
\cite[Theorem~2.2]{Mathias1973PrikryNative}.
\end{hbtheorem}

\begin{hbexample}{def:p06-tree-prikry-intersection-boundary}{Why all successor sets are not flattened into one intersection}
All sets \(\kappa\setminus\{\alpha\}\), \(\alpha<\kappa\), lie
in \(U\), but their intersection is empty. Thus a tree with
\(\kappa\) nodes cannot be flattened by intersecting all of its
successor sets merely using \(\kappa\)-completeness. The normal-free
proof uses separate grafts and fewer-than-\(\kappa\) intersections
\emph{at each retained node}. The normal comparison uses an explicitly
different diagonal argument. A principal ultrafilter would not even
contain all unbounded tails and is outside the standing hypothesis.
\end{hbexample}

\begin{proofboundary}{def:p06-tree-prikry-source-boundary}{The measure and forcing-family boundaries}
The tree construction uses the stated nonprincipal, internally complete
ground measure. The final dense comparison additionally uses normality.
One-element, supercompact, extender, Radin/Magidor and iteration
presentations have different condition systems and preservation arguments;
their conclusions do not follow merely from sharing the name Prikry.
No consistency implication is obtained by assuming a transitive ground
or a measurable cardinal from bare consistency.
\end{proofboundary}
